% 用法: xelatex SL_gap_nge2_exact_2n_switches_proof.tex (两遍, -output-directory=build)
% 主题: 相邻谱隙极值子的精确 2n 开关定理 (一维加权 Dirichlet, n >= 2)
% 来源: 用户提供 SL_gap_nge2_exact_2n_switches_proof_zh.pdf (2026-08), 本文为忠实转录,
%       不超出 PDF 内容, 保留原文"不声称首创"及文献优先级边界声明.
%       冻结英文证明: runs/R-20260809T125847Z-exact2n-switch/artifacts/full_proof.md,
%       SHA-256 9d4d8de8c18064717c8c7e2c913cdc1b914d086845a87cc55f34facdc0e0e5a9;
%       独立综合审计: routes/synthesis_audit/independent_audit.md,
%       SHA-256 3fe0d55cf4d09d3e0f0cc8ae23428bd69c058b6b01613ef0cb3ab008ddf3f918 (PASS).
%       Blueprint 绑定哈希: proposal 6ebe5137..., validation ae8bd1aa...,
%       review 290129f2... (approve), integration 70db8f75...,
%       合并后 blueprint 701e6793..., 合并后 evidence 1e83425e... (详见第 13 节).
% 审计 (2026-08-10, 项目侧): 解析逐条复核 PASS; 数值复跑全部通过 -
%       scripts/audit_nge2_pdfs.py Part A 40/40 (随机 bang-bang, R in {1.1..100}, n=2..5)
%       + Part B 16/16 (n1..n8 SUP/INF, R=4, 恰 2n 零点, q0>1, q1<-1, K+2D~0);
%       scripts/_hp_nge2.py (mpmath 50 位, n1_SUP/n4_INF/n8_INF);
%       scripts/_smooth_nge2.py (光滑振荡权, 零点公式与 W<0, 4/4).
% 文献: 未检索到直接等价已发表定理; Willner-Mahar 1979 为明确既有工作风险,
%       Sun 2022 只覆盖第一谱隙最小化 (INF 侧, n=1), 本文不声称首创 (见第 12 节).
\documentclass[UTF8]{ctexart}
\usepackage{amsmath, amssymb, amsthm}
\usepackage[margin=2.5cm]{geometry}
\usepackage[hyphens]{url}
\usepackage[hidelinks]{hyperref}
\usepackage{booktabs}
\usepackage{seqsplit}
\usepackage{graphicx}
\usepackage{longtable}
\emergencystretch 3em

\newtheorem{theorem}{定理}[section]
\newtheorem{proposition}[theorem]{命题}
\newtheorem{lemma}[theorem]{引理}
\newtheorem{corollary}[theorem]{推论}
\newtheorem{remark}[theorem]{注}
\newtheorem{definition}[theorem]{定义}

\title{相邻谱隙极值子的精确 $2n$ 开关定理\\ 一维加权 Dirichlet Sturm--Liouville 问题的完整证明}
\author{BVE research 项目 (Blueprint v2.2 数学研究证明包的中文自洽展开)}
\date{2026 年 8 月 (项目集成转录 2026-08-10)}

\begin{document}
\maketitle

\begin{abstract}
固定有限实数 $R>1$ 与整数 $n\ge2$, 考虑可测权重盒
\begin{equation*}
	A_R=\{\rho\in L^\infty(0,1):1\le\rho\le R\ \text{a.e.}\}
\end{equation*}
以及 Dirichlet 问题 $-u''=\lambda\rho u$、$u(0)=u(1)=0$. 本文证明: 相邻谱隙
$D_n(\rho)=\lambda_{n+1}(\rho)-\lambda_n(\rho)$ 的任一全局最大化子和任一全局最小化子
都恰有 $2n$ 个有效开关.

证明由四个相互锁定而不循环的部分组成. 第一, 固定加权归一化与端点定向后, 相邻特征函数的
Wronskian 在整个开区间严格为负, 故商 $Q=u_{n+1}/u_n$ 在 $u_n$ 的每个结点区间严格递减; 由此
得到切换函数 $F=\lambda_n u_n^2-\lambda_{n+1}u_{n+1}^2$ 的精确零点公式. 第二, Feynman--Hellmann
公式和盒约束的一侧局部变分给出最大、最小两种情形的完全饱和律; 所有简单零点都成为开关, 且不
存在额外开关. 第三, 在每个常密度块上构造两模态能量差 $K$, 其接口跳量等于密度跳量乘以 $F$,
因而在极值子上全局拼接为常数. 第四, 两次归一化贡献给出严格恒等式 $K=-2(\lambda_{n+1}-\lambda_n)$,
从而强迫两个端点斜率比满足 $q_0>1$ 与 $q_1<-1$, 精确零点数遂从可能的 $2n-2,2n-1,2n$ 收紧为
$2n$. 文末逐项给出定义、逻辑、边界与对抗审计; 数值计算只作为非证明性的压力测试. 本文不声称
文献首创性.
\end{abstract}

\tableofcontents

\section{问题设定、计数约定与主定理}
\subsection{可测权重盒与相邻谱隙}

全文固定一个有限实数 $R>1$. 允许的权重没有预设的连续性、分段性、对称性或单调性, 而是完整
的可测盒
\begin{equation}
	A_R=\left\{\rho\in L^\infty(0,1):1\le\rho(x)\le R\ \text{对几乎处处的}\ x\right\}.
	\label{eq:box}
\end{equation}
对每个 $\rho\in A_R$, 考虑弱意义下的加权 Dirichlet 特征值问题
\begin{equation}
	-u''=\lambda\rho u,\qquad u(0)=u(1)=0 .
	\label{eq:sl}
\end{equation}
它的特征值按严格递增次序记为
\begin{equation}
	0<\lambda_1(\rho)<\lambda_2(\rho)<\cdots,
\end{equation}
目标泛函为第 $n$ 个相邻谱隙
\begin{equation}
	D_n(\rho)=\lambda_{n+1}(\rho)-\lambda_n(\rho) .
	\label{eq:dn}
\end{equation}
附录 A 给出弱星紧性与谱连续性的自洽证明, 因此最大值与最小值都确实达到. 主证明从任取的一个
全局极值子出发; 因此所得结论自动覆盖极值集合中的每一个成员, 而不是仅覆盖某条数值分支.

\subsection{几乎处处等价与有效开关}

\begin{definition}[有限 bang--bang 代表元与有效开关]\label{def:switch}
称 $L^\infty$ 等价类 $\rho$ 具有一个有限 bang--bang 代表元, 如果存在分点
\begin{equation*}
	0=s_0<s_1<\cdots<s_m=1
\end{equation*}
与数值 $r_j\in\{1,R\}$, 使得 $\rho=r_j$ 几乎处处成立于每个正长度开区间 $(s_{j-1},s_j)$.
先合并所有相邻且取值相同的区间; 合并后每个内部公共端点 $s_1,\dots,s_{m-1}$ 称为一个有效开关.
有效开关数为 $m-1$. 固定的 Dirichlet 端点 $0,1$ 不计入开关.
\end{definition}

定义中的三个限定都不可省略: 零长度"块"不产生开关; 相邻同值的人为分点不产生开关; 在有限多个
接口点重新赋值不改变 $L^\infty$ 等价类, 也不改变谱或开关数.

\begin{theorem}[精确 $2n$ 有效开关]\label{thm:main}
对每个有限 $R>1$ 和每个整数 $n\ge2$, 泛函 $D_n$ 在 $A_R$ 上同时达到最大值与最小值.
进一步, $D_n$ 的每一个全局最大化子和每一个全局最小化子都具有有限 bang--bang 代表元, 并且该
代表元恰有 $2n$ 个有效内部开关.
\end{theorem}

这一定理不包含任何隐藏的"选支"条件: 极值子可以不唯一; 某个极值子与其反射可以不同; 证明也不
先验要求它们关于 $x=\tfrac12$ 对称.

\section{可测正权下的谱事实与端点定向}
\subsection{正则性、单性与加权归一化}

弱特征函数 $u\in H^1_0(0,1)$ 在一维中连续且有界. 因为 $u''=-\lambda\rho u\in L^\infty(0,1)$, 故
\begin{equation}
	u\in W^{2,\infty}(0,1)\subset C^1[0,1] .
	\label{eq:regularity}
\end{equation}
所有函数值和一阶导数的端点及接口迹因而都有经典意义. 对于 $L^\infty$ 系数, 一阶系统的积分方程
具有初值唯一性; 若非零解在某点同时满足 $u=u'=0$, 则它必恒为零. 于是每个零点都是简单零点.
两个属于同一特征值且满足左端 Dirichlet 条件的解由其左端导数唯一确定, 故每个特征值都是简单的.

后文只使用第 $n$ 与第 $n+1$ 个模态. 固定归一化和定向
\begin{equation}
	\int_0^1\rho u_k^2\,dx=1,\qquad u'_k(0)>0,\qquad k=n,n+1 .
	\label{eq:normalize}
\end{equation}
再记
\begin{equation}
	a=\lambda_n,\qquad b=\lambda_{n+1},\qquad D=b-a>0,\qquad
	c=\sqrt{\frac{a}{b}}\in(0,1),
	\label{eq:abc}
\end{equation}
以及切换函数
\begin{equation}
	F(x)=au_n(x)^2-bu_{n+1}(x)^2 .
	\label{eq:f}
\end{equation}

\subsection{结点数与严格交错}

下面回顾证明所需的振荡结论, 并说明它们为什么仍适用于仅可测的正权. 令 $s(x,\lambda)$ 解初值问题
\begin{equation*}
	s''+\lambda\rho s=0,\qquad s(0,\lambda)=0,\qquad s'(0,\lambda)=1 .
\end{equation*}
向量 $(s,s')$ 从不同时为零, 因此可连续选取 Pr\"ufer 角
\begin{equation*}
	s=r\sin\theta,\qquad s'=r\cos\theta,\qquad r>0,\qquad \theta(0,\lambda)=0 .
\end{equation*}
直接计算得, 几乎处处有
\begin{equation}
	\theta_x=\frac{s'^2+\lambda\rho s^2}{s^2+s'^2}
		=\cos^2\theta+\lambda\rho\sin^2\theta>0 .
	\label{eq:theta_x}
\end{equation}
因此每个零点对应 Pr\"ufer 角严格越过一个整数倍的 $\pi$. 对参数 $\lambda$ 作差商或对 Volterra
积分方程求导, 还得到端点角 $\theta(1,\lambda)$ 随 $\lambda$ 严格增加. 由此得到通常的结点定理:
属于 $\lambda_k$ 的特征函数在 $(0,1)$ 中恰有 $k-1$ 个简单零点.

\begin{lemma}[相邻模态严格交错]\label{lem:interlace}
令 $u=u_n$、$v=u_{n+1}$. 若 $\alpha_1<\cdots<\alpha_{n-1}$ 是 $u$ 的内部零点,
$\beta_1<\cdots<\beta_n$ 是 $v$ 的内部零点, 则
\begin{equation}
	0<\beta_1<\alpha_1<\beta_2<\cdots<\alpha_{n-1}<\beta_n<1 .
	\label{eq:interlace}
\end{equation}
\end{lemma}

\begin{proof}
设 $(\alpha,\tilde\alpha)$ 是 $u$ 的任一结点区间, 允许其一为 Dirichlet 端点. 反设 $v$ 在该区间内
无零点. 改变符号后可令区间内 $u,v>0$, 并定义
\begin{equation*}
	W=v'u-vu' .
\end{equation*}
由两个方程得
\begin{equation}
	W'=-(b-a)\rho uv<0\quad\text{a.e.}
	\label{eq:wprime}
\end{equation}
另一方面, $u'(\alpha)>0$、$u'(\tilde\alpha)<0$, 而连续延拓的 $v$ 在两端非负, 故
\begin{equation*}
	W(\alpha)=-v(\alpha)u'(\alpha)\le0,\qquad
	W(\tilde\alpha)=-v(\tilde\alpha)u'(\tilde\alpha)\ge0,
\end{equation*}
与 $W$ 严格下降矛盾. 因此 $u$ 的每个结点区间至少含有一个 $v$ 零点. $u$ 有 $n$ 个结点区间, 而
$v$ 恰有 $n$ 个内部零点, 所以每个区间恰含一个, 且没有共同内部零点, 得到 \eqref{eq:interlace}.
\end{proof}

\subsection{右端导数的奇偶符号}

由 \eqref{eq:normalize}, 每个 $u_k$ 在左端附近为正. 它经过 $k-1$ 个简单内部零点后, 在最后一个
结点区间的符号为 $(-1)^{k-1}$. 而 $u_k(x)=u'_k(1)(x-1)+o(1-x)$, 故
\begin{equation}
	\operatorname{sgn}u'_k(1)=(-1)^k .
	\label{eq:parity}
\end{equation}
定义两个端点斜率比
\begin{equation}
	q_0=\frac{u'_{n+1}(0)}{u'_n(0)},\qquad q_1=\frac{u'_{n+1}(1)}{u'_n(1)},\qquad
	q_0>0,\qquad q_1<0 .
	\label{eq:q}
\end{equation}
在尚未使用极值性时, 已经知道
\begin{equation*}
	q_0>0,\qquad q_1<0 .
\end{equation*}

\section{Wronskian 在整个开区间严格为负}

继续令 $u=u_n$、$v=u_{n+1}$, 并固定本文的符号约定
\begin{equation}
	W=v'u-vu' .
	\label{eq:wdef}
\end{equation}
由特征方程, 几乎处处有
\begin{equation}
	W'=v''u-vu''=-(b-a)\rho uv=-D\rho uv .
	\label{eq:wprime2}
\end{equation}

\begin{lemma}[全局严格 Wronskian 符号]\label{lem:wneg}
在上述归一化和符号约定下,
\begin{equation}
	W(x)<0\qquad(0<x<1) .
	\label{eq:wneg}
\end{equation}
\end{lemma}

\begin{proof}
在 $v$ 的第 $j$ 个零点 $\beta_j$ 处, 交错次序与从左端开始的符号交替给出
\begin{equation*}
	\operatorname{sgn}v'(\beta_j)=(-1)^j,\qquad
	\operatorname{sgn}u(\beta_j)=(-1)^{j-1},
\end{equation*}
所以
\begin{equation}
	W(\beta_j)=v'(\beta_j)u(\beta_j)<0 .
	\label{eq:wbeta}
\end{equation}
在 $u$ 的第 $j$ 个内部零点 $\alpha_j$ 处,
\begin{equation*}
	\operatorname{sgn}u'(\alpha_j)=(-1)^j,\qquad
	\operatorname{sgn}v(\alpha_j)=(-1)^j,
\end{equation*}
故
\begin{equation}
	W(\alpha_j)=-v(\alpha_j)u'(\alpha_j)<0 .
	\label{eq:walpha}
\end{equation}
在交错序列 \eqref{eq:interlace} 的任意两个相邻点之间, $uv$ 符号固定, 故 \eqref{eq:wprime2}
表明 $W$ 单调; 该小区间的两个端值由 \eqref{eq:wbeta}--\eqref{eq:walpha} 都为负, 因此内部也为
负. 在首区间 $(0,\beta_1)$ 上 $uv>0$, 所以 $W$ 从 $W(0)=0$ 严格下降. 在尾区间 $(\beta_n,1)$ 上
$uv<0$, 所以 $W$ 从负值严格上升到 $W(1)=0$. 所有区间合并即得 \eqref{eq:wneg}.
\end{proof}

这里的严格性至关重要. 若只知道 $W\le0$, 商函数可能出现平台, 从而无法排除切换函数的重零点或
零区间. 引理 \ref{lem:wneg} 完全由振荡与交错推出, 尚未使用极值性或开关数信息.

\section{切换函数的精确零点公式}

令
\begin{equation*}
	0=x_0<x_1<\cdots<x_{n-1}<x_n=1
\end{equation*}
为 $u_n$ 的全部零点 (含端点), 并在每个结点区间 $I_j=(x_{j-1},x_j)$ 上定义
\begin{equation}
	Q=\frac{u_{n+1}}{u_n} .
	\label{eq:qdef}
\end{equation}
由引理 \ref{lem:wneg},
\begin{equation}
	Q'=\frac{u'_{n+1}u_n-u_{n+1}u'_n}{u_n^2}=\frac{W}{u_n^2}<0 .
	\label{eq:qprime}
\end{equation}
因此 $Q$ 在每个结点区间严格递减.

\subsection{每个结点区间的值域}

由严格交错和简单零点, $Q$ 的单侧极限为
\begin{center}
\begin{tabular}{lcc}
\toprule
结点区间 & 左端极限 & 右端极限\\
\midrule
首区间 $I_1$ & $q_0$ & $-\infty$\\
中间区间 $I_j$, $2\le j\le n-1$ & $+\infty$ & $-\infty$\\
尾区间 $I_n$ & $+\infty$ & $q_1$\\
\bottomrule
\end{tabular}
\end{center}
首尾有限极限来自洛必达型的简单零点展开, 例如 $\lim_{x\downarrow0}Q(x)=u'_{n+1}(0)/u'_n(0)=q_0$.

在 $u_n\ne0$ 处, 利用 \eqref{eq:abc} 有恒等式
\begin{equation}
	F=au_n^2-bu_{n+1}^2=bu_n^2\left(c^2-Q^2\right) .
	\label{eq:fq}
\end{equation}
而在 $u_n$ 的内部零点处, 严格交错保证 $u_{n+1}\ne0$, 故 $F<0$. 所以 $F$ 的全部内部零点恰是各
结点区间内的横截 $Q=+c$ 或 $Q=-c$.

\begin{proposition}[精确零点计数]\label{prop:count}
设 $Z(F;(0,1))=\{x\in(0,1):F(x)=0\}$. 则
\begin{equation}
	\#Z(F;(0,1))=2n-2+\mathbf 1_{\{q_0>c\}}+\mathbf 1_{\{q_1<-c\}} .
	\label{eq:count}
\end{equation}
而且这些零点全部简单, 因而每个零点都伴随 $F$ 的符号改变.
\end{proposition}

\begin{proof}
首区间中, $Q$ 从正数 $q_0$ 严格降至 $-\infty$, 所以必穿过 $-c$ 一次; 它穿过 $+c$ 当且仅当
$q_0>c$. 每个中间区间中, $Q$ 从 $+\infty$ 严格降至 $-\infty$, 故恰穿过 $+c$ 与 $-c$ 各一次.
尾区间中, $Q$ 从 $+\infty$ 严格降至负数 $q_1$, 所以必穿过 $+c$ 一次; 它穿过 $-c$ 当且仅当
$q_1<-c$. 总数为
\begin{equation*}
	1+2(n-2)+1+\mathbf 1_{\{q_0>c\}}+\mathbf 1_{\{q_1<-c\}},
\end{equation*}
即 \eqref{eq:count}.

若 $q_0=c$, 等号只发生在被排除的端点 $x=0$, 严格递减使区间内立即有 $Q<c$, 所以不新增内部
零点. 若 $q_1=-c$, 同理等号只发生在 $x=1$. 这解释了 \eqref{eq:count} 中必须使用严格指标.

最后, 在 $F=0$ 的内部点, $u_n\ne0$、$Q=\pm c$. 微分 \eqref{eq:fq} 得
\begin{equation}
	F'=-2bu_n^2QQ'\ne0,
	\label{eq:fprime}
\end{equation}
因为 $c>0$ 且 $Q'<0$. 故所有零点简单并变号.
\end{proof}

\begin{remark}[最低编号 $n=2$]\label{rem:n2}
当 $n=2$ 时没有中间结点区间, 上述求和中的 $2(n-2)$ 为零; 首尾各有一个必然横截, 再加两个端点
指标. 因此公式仍原样成立, 没有暗中使用 $n\ge3$.

在一般正权下, 两个端点指标未必都为 $1$, 所以 \eqref{eq:count} 只给出 $2n-2,2n-1,2n$ 三种可能.
真正把三分情形压缩为精确 $2n$ 的, 是第 6 节的极值子能量不变量.
\end{remark}

\section{Feynman--Hellmann 公式与盒约束的完全饱和}
\subsection{任意有界方向的一阶公式}

\begin{lemma}[Feynman--Hellmann]\label{lem:fh}
令 $h\in L^\infty(0,1)$, 并令 $u_k$ 按 \eqref{eq:normalize} 归一化. 则简单特征值沿
$\rho+\varepsilon h$ 的方向导数为
\begin{equation}
	\lambda'_k[h]=-\lambda_k\int_0^1 hu_k^2\,dx .
	\label{eq:fh}
\end{equation}
因此
\begin{equation}
	D'_n[h]=\int_0^1 hF\,dx .
	\label{eq:fhgap}
\end{equation}
\end{lemma}

\begin{proof}
在 $H^1_0(0,1)$ 上以 $\langle f,g\rangle_H=\int f'g'$ 为内积, 定义紧自伴算子 $T_\rho$ 使
\begin{equation*}
	\langle T_\rho f,g\rangle_H=\int_0^1\rho fg\,dx .
\end{equation*}
有精确的仿射扰动 $T_{\rho+\varepsilon h}=T_\rho+\varepsilon T_h$. 简单孤立特征值的标准紧自伴扰动论
保证对应特征值与定向特征向量可微. 也可直接对弱特征方程作差商并利用简单性取极限.

形式上对
\begin{equation*}
	\int u'_{k,\varepsilon}\phi'=\lambda_{k,\varepsilon}\int(\rho+\varepsilon h)u_{k,\varepsilon}\phi
\end{equation*}
在 $\varepsilon=0$ 求导, 再取 $\phi=u_k$. 由原方程测试 $\dot u_k$ 后, 含 $\dot u_k$ 的两项完全
抵消, 留下
\begin{equation*}
	0=\lambda'_k\int\rho u_k^2+\lambda_k\int hu_k^2 .
\end{equation*}
使用加权单位归一化即得 \eqref{eq:fh}. 对 $k=n+1$ 与 $k=n$ 两式相减便得 \eqref{eq:fhgap}.
\end{proof}

\subsection{最大化与最小化的两种饱和律}

\begin{proposition}[完全盒饱和]\label{prop:saturation}
设 $\rho^*$ 为一个全局极值子, 并由它的归一化特征函数定义 $F$. 若 $\rho^*$ 是全局最大化子, 则
\begin{equation}
	\rho^*=R\ \text{a.e. 于}\ \{F>0\},\qquad
	\rho^*=1\ \text{a.e. 于}\ \{F<0\} .
	\label{eq:satmax}
\end{equation}
若 $\rho^*$ 是全局最小化子, 则
\begin{equation}
	\rho^*=1\ \text{a.e. 于}\ \{F>0\},\qquad
	\rho^*=R\ \text{a.e. 于}\ \{F<0\} .
	\label{eq:satmin}
\end{equation}
\end{proposition}

\begin{proof}
先看最大化. 对 $\delta>0$, 置
\begin{equation*}
	A_\delta=\{\rho^*\le R-\delta\},\qquad B_\delta=\{\rho^*\ge1+\delta\}.
\end{equation*}
若 $\phi\ge0$ 有界且支撑在 $A_\delta$, 则对充分小的 $\varepsilon>0$, 方向 $h=\phi$ 可行. 最大性
与 \eqref{eq:fhgap} 给出 $\int\phi F\le0$, 从而
\begin{equation}
	F\le0\quad\text{a.e. 于}\ \{\rho^*<R\}.
	\label{eq:sat1}
\end{equation}
同理, 在 $B_\delta$ 上取 $h=-\phi$, 得到
\begin{equation}
	F\ge0\quad\text{a.e. 于}\ \{\rho^*>1\}.
	\label{eq:sat2}
\end{equation}
若 $F>0$, 式 \eqref{eq:sat1} 排除 $\rho^*<R$, 故 $\rho^*=R$; 若 $F<0$, 式 \eqref{eq:sat2} 排除
$\rho^*>1$, 故 $\rho^*=1$. 这就是 \eqref{eq:satmax}.

最小化时所有可行一侧方向满足 $D'_n[h]\ge0$. 完全相同的局部测试把上述符号反转, 得到
\eqref{eq:satmin}.
\end{proof}

\begin{center}
\begin{tabular}{lcc}
\toprule
 & $F>0$ & $F<0$\\
\midrule
全局最大化子 & $\rho=R$ & $\rho=1$\\
全局最小化子 & $\rho=1$ & $\rho=R$\\
\bottomrule
\end{tabular}
\end{center}

\subsection{零点集合与有效开关集合双向相等}

\begin{proposition}[零点就是开关]\label{prop:switchzero}
对任一全局最大化子或全局最小化子, 存在有限 bang--bang 代表元, 并且
\begin{equation}
	\{\text{有效开关}\}=Z(F;(0,1)) .
	\label{eq:switchzero}
\end{equation}
\end{proposition}

\begin{proof}
命题 \ref{prop:count} 已在使用极值性之前证明: $F$ 的内部零点有限且全部简单. 因此
$(0,1)\setminus Z(F;(0,1))$ 只有有限个连通分支, 且 $F$ 在每个分支上保持严格定号. 饱和律
\eqref{eq:satmax} 或 \eqref{eq:satmin} 随即把 $\rho^*$ 在每个分支上的几乎处处值固定为 $1$ 或
$R$, 故得到有限 bang--bang 代表元.

每个 $F$ 零点都简单, 所以穿过它时 $F$ 变号. 无论采用最大化还是最小化的材料指派, 变号都会使
相邻两个分支取不同的盒端点, 因而每个零点都是有效开关. 反过来, 在同一个定号分支内, 饱和律只
允许一个几乎处处常值; 合并相邻同值块并忽略接口点赋值后, 不可能还有额外有效开关. 两方向合并
即得 \eqref{eq:switchzero}.
\end{proof}

这一步已经建立有限分块, 因而下一节在常密度块上定义能量并非预设极值子分段, 而是从变分饱和
严格推出的结果.

\section{全局块能量不变量}

固定任一全局极值子, 并采用命题 \ref{prop:switchzero} 给出的有限步函数代表元. 在密度为常数
$r\in\{1,R\}$ 的一个开块上, 定义
\begin{equation}
	K(x)=\left(u'_n(x)^2+ar\,u_n(x)^2\right)
		-\left(u'_{n+1}(x)^2+br\,u_{n+1}(x)^2\right) .
	\label{eq:kdef}
\end{equation}

\begin{proposition}[能量差的全局拼接]\label{prop:k}
式 \eqref{eq:kdef} 在所有常密度块上拼接成同一个全局常数, 并且
\begin{equation}
	K=-2(b-a)=-2D<0 .
	\label{eq:kval}
\end{equation}
\end{proposition}

\begin{proof}
首先, 在常密度块内, 对任一 $k=n,n+1$ 有
\begin{equation}
	\frac{d}{dx}\left(u_k'^2+\lambda_k r u_k^2\right)
		=2u'_ku''_k+\lambda_k r u_k^2=0\quad\text{a.e.}
	\label{eq:blockconst}
\end{equation}
因此 \eqref{eq:kdef} 在每个块内为常数.

设 $s$ 是一个有效开关, 其本质单侧密度值为 $r_-$、$r_+$. 正则性 \eqref{eq:regularity} 保证两个
特征函数及其一阶导数在 $s$ 处有共同连续迹, 所以动能项无跳跃; 唯一可能的跳跃来自密度值. 精确
计算为
\begin{equation}
	K(s_+)-K(s_-)=(r_+-r_-)\left(au_n(s)^2-bu_{n+1}(s)^2\right)
		=(r_+-r_-)F(s)=0 .
	\label{eq:jump}
\end{equation}
最后一个等号使用 \eqref{eq:switchzero}: 每个有效开关正是 $F$ 的零点. 故相邻块的常数完全相同,
所有块拼成一个全局的几乎处处常数. 注意 $\rho(s)$ 在单点的具体赋值从未进入 \eqref{eq:jump};
只使用本质单侧值和 $C^1$ 迹.

现在在长度为 $1$ 的区间上积分这个常数. 弱特征方程以 $u_k$ 自测, 并使用 \eqref{eq:normalize},
得到
\begin{equation}
	\int_0^1u_n'^2\,dx=a,\qquad \int_0^1u_{n+1}'^2\,dx=b .
	\label{eq:kinetic}
\end{equation}
又有
\begin{equation}
	\int_0^1\rho u_n^2\,dx=a,\qquad \int_0^1\rho u_{n+1}^2\,dx=b .
	\label{eq:potential}
\end{equation}
于是
\begin{equation}
	K=\int_0^1K(x)\,dx
		=\underbrace{(a-b)}_{\text{动能差}}+\underbrace{(a-b)}_{\text{加权势能差}}
		=-2(b-a) .
	\label{eq:kintegral}
\end{equation}
这同时解释了不可遗漏的因子 $2$: 它来自两个彼此独立、数值相同的归一化贡献.
\end{proof}

\begin{remark}[为什么任意 bang--bang 权重不够]\label{rem:bangbang}
若一个非极值 bang--bang 权重在某接口 $s$ 满足 $F(s)\ne0$, 则 \eqref{eq:jump} 给出非零跳量,
块能量无法全局拼接. 故端点刚性不是任意分块权重的性质, 而是"完全饱和使所有接口落在
$F=0$"的极值结构后果.
\end{remark}

\section{端点刚性与精确 \texorpdfstring{$2n$}{2n} 结论}

因为 $K$ 的所有块常数相同, 其首尾单侧迹都等于 \eqref{eq:kval}. Dirichlet 条件使端点处的势能项
消失, 故
\begin{equation}
	u'_n(0)^2-u'_{n+1}(0)^2=-2D,\qquad
	u'_n(1)^2-u'_{n+1}(1)^2=-2D .
	\label{eq:endpoint}
\end{equation}
由于 $D>0$, 连续的高一阶模态在两个端点都有严格更大的斜率绝对值:
\begin{equation}
	|u'_{n+1}(0)|>|u'_n(0)|,\qquad |u'_{n+1}(1)|>|u'_n(1)| .
	\label{eq:abs}
\end{equation}
左端定向 \eqref{eq:normalize} 与右端奇偶式 \eqref{eq:parity} 把绝对值结论升级为有向比值结论
\begin{equation}
	q_0>1,\qquad q_1<-1 .
	\label{eq:qstrict}
\end{equation}
另一方面 $0<c<1$, 所以
\begin{equation*}
	q_0>1>c,\qquad q_1<-1<-c .
\end{equation*}
命题 \ref{prop:count} 中两个严格指标都等于 $1$, 从而
\begin{equation}
	\#Z(F;(0,1))=2n .
	\label{eq:finalcount}
\end{equation}
再由命题 \ref{prop:switchzero}, 有效开关数就是这个零点数. 因此每个任取的最大化子或最小化子都
恰有 $2n$ 个有效开关, 定理 \ref{thm:main} 得证. \qed

\begin{corollary}[两种极值子的材料起止顺序]\label{cor:materials}
合并相邻同值块后, 每个极值子恰有 $2n+1$ 个正长度块, 并严格交替取值. 最大化子的首块和尾块都
取 $1$; 最小化子的首块和尾块都取 $R$.
\end{corollary}

\begin{proof}
由 \eqref{eq:qstrict} 与 \eqref{eq:fq}, $F$ 在两个端点内侧都为负. 全部 $2n$ 个零点简单, 故符号
逐次交替; $2n$ 为偶数, 所以首尾符号相同. 最后应用表 \eqref{eq:satmax}--\eqref{eq:satmin} 的
材料指派即可.
\end{proof}

推论 \ref{cor:materials} 确定了材料的交替次序, 但没有确定开关位置、块长、对称性或唯一性.

\section{逻辑依赖、前提契约与非循环性审计}

为避免把结论偷偷作为前提, 下面按实际使用顺序列出证明链.

\begin{center}
\resizebox{\textwidth}{!}{%
\begin{tabular}{lll}
\toprule
步骤 & 输入 & 严格输出\\
\midrule
1 & $1\le\rho\le R<\infty$, 一维 Dirichlet 方程 &
$W^{2,\infty}$ 正则性、特征值单性、精确结点数与相邻严格交错.\\
2 & 交错、端点定向、$W'=-(b-a)\rho uv$ &
$W<0$ 与每个结点区间的商极限.\\
3 & 简单特征值扰动与盒约束的一侧可行方向 &
精确零点式 \eqref{eq:count}, 且所有零点简单; 此时尚未使用极值性.\\
4 & 最大、最小两套完全饱和律.\\
5 & 已得到的有限分块、接口处 $F=0$ &
有限简单零点与完全饱和.\\
6 & 有限 bang--bang 代表元, 并且开关集合恰等于零点集合.\\
7 & 两个加权单位归一化与自测恒等式 &
块能量在接口无跳, 拼成全局常数.\\
8 & Dirichlet 端点、定向与奇偶符号 &
$K=-2(b-a)<0$.\\
9 & 精确零点式与开关--零点等式 &
$q_0>1$、$q_1<-1$, 两个端点指标都为 $1$.\\
\bottomrule
\end{tabular}}
\end{center}

\subsection{非循环性要点}

能量不变量的构造只需要"零点有限且简单"以及"每个开关都是零点", 并不需要预先知道零点总数为
$2n$. 具体顺序是:
\begin{equation*}
	\text{Wronskian 严格符号}\Rightarrow\text{零点有限简单}\Rightarrow
	\text{饱和与有限分块}\Rightarrow\text{能量拼接}\Rightarrow
	\text{端点指标}\Rightarrow 2n .
\end{equation*}
所以最后一步的精确计数从未被用于建立倒数第二步的端点刚性.

\subsection{量词与前提边界}

\begin{itemize}
	\item $R$ 是任意但固定的有限实数, 且严格满足 $R>1$.
	\item $n$ 是任意整数 $n\ge2$, 包括没有中间结点区间的 $n=2$.
	\item 证明从任意一个全局极值子出发, 故结论量词是"每一个", 不是"至少存在一个".
	\item 最大化与最小化只在饱和指派方向上不同; 其余零点与能量论证完全共用.
	\item 既有的"开关数至多 $2n$"定理与本结论一致, 但本文没有用它制造缺失的下界. 真正的下界
		来自精确零点式和端点能量刚性.
\end{itemize}

\section{边界情形与对抗性审计}
\subsection{定义与边界检查}

\begin{center}
\resizebox{\textwidth}{!}{%
\begin{tabular}{ll}
\toprule
检查对象 & 处理结果\\
\midrule
$n=2$ &
中间结点区间族为空; 首尾两个必然横截加两个端点指标,\\
& 能量论证给出四个开关.\\
$R\downarrow1$ &
对每个固定 $R>1$ 论证都严格成立, 不需要 $R-1$ 的统一正下界.\\
& $R=1$ 时允许集退化且开关为零, 因此它被定理明确排除.\\
任意大的有限 $R$ &
证明没有大对比近似; 只要 $R<\infty$, 正则性、变分和能量恒等式均保持.\\
& $R=\infty$ 不在当前允许集内.\\
端点等号 $q_0=c$ 或 $q_1=-c$ &
原始零点公式用严格指标, 等号只发生在不计数的端点; 能量恒等式随后证明\\
& 更强的 $q_0>1>c$、$q_1<-1<-c$, 故极值子不会落在等号情形.\\
几乎处处代表元 &
只使用块的本质单侧密度值与特征函数的共同 $C^1$ 迹. 改变有限个接口点\\
& 的赋值不影响任何结论.\\
Dirichlet 端点 &
端点用于读取全局能量常数, 但根据定义不计为开关.\\
最大化与最小化 &
饱和指派反转; 简单零点在两种指派下都使材料改变, 接口能量跳量都为零.\\
非唯一与反射 &
未选择对称分支, 也未用唯一性; 任一极值子及其可能不同的反射都分别满足\\
& 同一证明.\\
端点定向 &
明确固定 $u'_k(0)>0$. 若不固定特征函数符号, $q_0,q_1$ 的有向陈述没有意义;\\
& $F$ 本身则不受符号翻转影响.\\
奇异弧或正测度零集 &
精确零点命题给出有限且简单的零集, 排除了 $F=0$ 的正测度区间, 也排除了\\
& 盒内部值占据正测度的奇异弧.\\
\bottomrule
\end{tabular}}
\end{center}

\subsection{针对潜在错误的攻击}

\begin{enumerate}
	\item \emph{Wronskian 约定反号}. 本文始终使用 $W=u'_{n+1}u_n-u_{n+1}u'_n$, 因此 $W<0$
		且 $Q'<0$. 若换用相反定义, 后续所有符号必须同时反转, 不能混用.
	\item \emph{"开关推出零点"并不自动给出"零点推出开关"}. 第二方向依赖两个事实: 完全盒饱和,
		以及每个 $F$ 零点简单并变号. 命题 \ref{prop:switchzero} 明确证明了两方向.
	\item \emph{接口处遗漏密度跳量}. 正确公式是 $K(s_+)-K(s_-)=(r_+-r_-)F(s)$, 不是无条件为零.
		只有极值饱和把接口放在 $F=0$ 上时, 跳量才消失.
	\item \emph{漏掉 $K=-2D$ 的因子 $2$}. 式 \eqref{eq:kintegral} 分别列出动能差与加权势能差;
		二者各贡献一次 $a-b$.
	\item \emph{右端只得到绝对值}. 能量先给出 $|u'_{n+1}(1)|>|u'_n(1)|$, 再由固定左端定向后的
		奇偶式 \eqref{eq:parity} 得到负比值, 故是 $q_1<-1$, 不是仅有 $|q_1|>1$.
	\item \emph{把任意 bang--bang 测试权重当成极值子}. 说明 \ref{rem:bangbang} 表明一般接口可有
		非零能量跳量; 因此一般权重的端点斜率比失败不构成反例.
	\item \emph{把数值拓扑搜索当作完备性证明}. 有限网格、有限开关参数化或局部驻点搜索都不能覆盖
		完整的 $L^\infty$ 盒. 本文的普遍量词只由解析证明承担.
\end{enumerate}

独立综合审计重新推导了正则性、交错、Wronskian、精确零点式、两套一侧变分、零点与开关双向等同、
接口拼接、因子 $2$ 与端点奇偶, 结论为 PASS, 且无未解除证明义务. 该审计是证明包的独立复核,
不代替尚待完成的正式 Blueprint proposal review 与确定性接收.

\section{数值压力测试的证据边界}

独立计算路线对解析恒等式进行了高精度压力测试, 覆盖 $n=2,3$ 以及从接近 $1$ 到较大有限值的
多种对比度. 测试内容包括:
\begin{itemize}
	\item 在完全饱和的驻点分支上检查每个块内的能量常数、接口跳量和全局恒等式 $K=-2D$;
	\item 对随机正权检查商函数单调性与精确零点公式 \eqref{eq:count};
	\item 搜索较少开关的受限拓扑分支, 并在完整区间上复查饱和符号, 而不只检查参数化族内部的
		驻点条件;
	\item 在近简并或高对比情形用高精度重新计算, 以识别双精度造成的伪接口跳跃.
\end{itemize}
这些实验没有发现解析恒等式的反例, 但其逻辑地位严格限于对抗性核查: 它们不证明极值存在, 不证明
候选点全局最优, 不证明拓扑搜索完备, 也不承担"对所有 $R,n$ 与每个极值子"的量词. 主定理的每一步
均可在删除全部数值材料后独立成立.

\section{已经解决的范围与仍开放的问题}
\subsection{本证明已经闭合的内容}

在当前一维 Dirichlet、正可测盒权与有限 $R>1$ 的设定中, 已经证明:
\begin{itemize}
	\item 最大值和最小值都达到;
	\item 每个全局极值子都有限 bang--bang;
	\item 每个极值子的有效开关数不仅至多为 $2n$, 而且必定恰等于 $2n$;
	\item 每个极值子恰有 $2n+1$ 个交替材料块; 最大化子从 $1$ 开始并以 $1$ 结束, 最小化子从
		$R$ 开始并以 $R$ 结束;
	\item 这些结论不依赖对称性、唯一性或选定数值分支.
\end{itemize}

\subsection{仍然开放或未在本文解决的问题}

\begin{enumerate}
	\item \emph{开关位置与块长}. 给定 $(n,R)$ 后, $2n$ 个开关的精确位置是否由有限显式方程唯一
		确定? 是否存在可认证的闭式或半解析构造?
	\item \emph{反射对称性}. 是否至少存在一个关于 $x=\tfrac12$ 对称的最大化子或最小化子? 更强地,
		每个极值子是否都对称?
	\item \emph{唯一性与完整分类}. 除去反射后极值子是否唯一? 若不唯一, 全部全局极值子的连通分支、
		分岔点与等值机制如何分类?
	\item \emph{最优值}. $\max D_n$ 与 $\min D_n$ 是否存在闭式公式、锐界或高效且可严格认证的
		有限维算法?
	\item \emph{参数与渐近}. 当 $R\downarrow1$、$R\to\infty$ 或 $n\to\infty$ 时, 开关位置和最优
		谱隙有哪些统一渐近规律? 极限与优化能否交换?
	\item \emph{稳定性}. 偏离饱和开关位置时, 谱隙损失是否满足二次强制估计? 这关系到数值算法的
		条件数与近极值子的结构稳定性.
	\item \emph{模型推广}. Robin 或混合边界、一般系数 Sturm--Liouville 算子、非相邻谱隙、积分或
		体积分数约束下, 精确开关计数如何改变?
	\item \emph{文献优先级}. 当前检索尚不足以判定精确 $2n$ 全量词定理与早期一般特征值泛函极值
		理论之间的包含关系; 正式发表前仍需逐页核对早期文献.
\end{enumerate}

\section{文献定位与首创性边界}

本文所用的一维 Sturm--Liouville 振荡、简单谱与 Pr\"ufer 方法属于标准理论; 可参见 Teschl 的常
微分方程教材以及 Kong--Zettl 对正则 Sturm--Liouville 特征值问题的研究. Willner--Mahar 研究了一
般特征值函数的极值刻画, Sun 研究了振动弦第一谱隙的最小化问题. 这些来源说明本问题处在既有极值
谱理论的连续传统中.

截至本证明包形成时, 有限的主来源检索没有核实到一个以完全相同范围陈述的定理: 所有 $n\ge2$、完整
可测盒、最大与最小两端、每一个极值子、以及精确 $2n$ 个有效开关同时出现. 但"尚未核实到"不等于
"文献中不存在", 更不构成首创性或优先权证明. 因此本文只陈述数学推导在当前证明包中已经闭合, 明确
把 novelty/priority 标记为未知, 不使用"首次"、"全新"等措辞.

\section{Blueprint 证明包绑定与待接收字段}
\subsection{已经冻结并独立核验的数学对象}

本文展开的英文冻结证明为
\texttt{\seqsplit{runs/R-20260809T125847Z-exact2n-switch/artifacts/full\_proof.md}}, 其 SHA--256 为
\begin{center}
	\texttt{\seqsplit{sha256:9d4d8de8c18064717c8c7e2c913cdc1b914d086845a87cc55f34facdc0e0e5a9}}.
\end{center}
独立综合审计文件为
\texttt{\seqsplit{runs/R-20260809T125847Z-exact2n-switch/routes/synthesis\_audit/independent\_audit.md}},
其 SHA--256 为
\begin{center}
	\texttt{\seqsplit{sha256:3fe0d55cf4d09d3e0f0cc8ae23428bd69c058b6b01613ef0cb3ab008ddf3f918}}.
\end{center}
该审计结论为 PASS. 定义、逻辑、边界与对抗四份作者侧审计也分别绑定到冻结哈希; 它们与独立综合
审计都不修改 canonical Blueprint.

研究本轮启动并完成初始问题图接收后, 供独立审计只读比对的 canonical 快照为:
\begin{itemize}
	\item 启动快照 blueprint: \texttt{\seqsplit{sha256:295b5663c5466ed0f8d72e571d5df84630ffa9f204c01a5246e786b9269b763}};
	\item 启动快照 evidence inventory: \texttt{\seqsplit{sha256:02529f9374d833dab60c7f0650cc0b07431f2d172195100b069de0e7759107be}}.
\end{itemize}
这些是接收新定理之前的快照, 不得误写成最终合并后状态.

\subsection{最终 proposal、review 与 integration 接收记录}

下列字段均来自已冻结的 R2 proposal、确定性 validation、独立正式 review、确定性 integration
receipt 以及合并后的只读闭合查询; 没有任何字段来自推测或手工改写 canonical 文件.
\begin{center}
\resizebox{\textwidth}{!}{%
\begin{tabular}{ll}
\toprule
已核验的 Blueprint 接收字段 & 值\\
\midrule
最终 proposal SHA--256 & \texttt{6ebe5137af517358c1bf5f5f6ccb2fa129183f6a9bc9bf84cb4faa14de7a623e}\\
最终 validation SHA--256 & \texttt{ae8bd1aa4da1c2c6fa8512cee19d8c57124cae7089b4f8f16b39aa4ed2f62041}\\
独立正式 review 决定 & approve (findings 与 required actions 均为空)\\
独立正式 review SHA--256 & \texttt{290129f244441ecde4b530824dc59e27322254334548c65299caf36fa64252af}\\
确定性 integration receipt SHA--256 & \texttt{70db8f756e3193ee036d2c4a3da6362945d0348756ada81c19947f9cdf0edcfc}\\
合并后 blueprint SHA--256 & \texttt{701e67939345ee2db64c873f580ea8f2fb09f415664b0a29d5827fb450d6111f}\\
合并后 evidence inventory SHA--256 & \texttt{1e83425e2987078a6fb26f14b37fdef6ce26067798abd3a96683a490ff3b4887}\\
合并后 closure/frontier 状态 & PASS: 目标已进入可信闭包, 研究 frontier 为空\\
\bottomrule
\end{tabular}}
\end{center}
数学定理的证明真值不依赖这些流程哈希; 但"该结果已被 Blueprint v2.2 canonical 图正式接收"是另
一个独立命题, 只有 proposal 验证、独立 review、确定性 integration 与 post-merge closure 查询
全部成功后才能陈述.

\appendix
\section{极值达到性的自洽证明}

本附录补足主定理中最大值与最小值确实达到这一上游事实. 令 $H=H^1_0(0,1)$, 以内积
\begin{equation*}
	\langle f,g\rangle_H=\int_0^1 f'g'\,dx
\end{equation*}
赋予 Hilbert 结构. 对 $\rho\in A_R$, 由 Riesz 表示定义 $T_\rho:H\to H$:
\begin{equation}
	\langle T_\rho f,g\rangle_H=\int_0^1\rho fg\,dx .
	\label{eq:ta}
\end{equation}
一维紧嵌入 $H\hookrightarrow C[0,1]$ 表明 $T_\rho$ 是紧、正、自伴算子. 其正特征值
$\mu_1\ge\mu_2\ge\cdots\downarrow0$ 与 \eqref{eq:sl} 的特征值满足
\begin{equation}
	\mu_k(\rho)=\lambda_k(\rho)^{-1} .
	\label{eq:muk}
\end{equation}

\begin{lemma}[弱星收敛推出算子范数收敛]\label{lem:opnorm}
若 $\rho_j,\rho\in A_R$ 且 $\rho_j\overset{*}{\rightharpoonup}\rho$ 于 $L^\infty(0,1)$, 则
\begin{equation*}
	\|T_{\rho_j}-T_\rho\|_{H\to H}\to0 .
\end{equation*}
\end{lemma}

\begin{proof}
由 \eqref{eq:ta}, 所求范数等于
\begin{equation*}
	\sup_{\|f\|_H\le1,\|g\|_H\le1}\left|\int_0^1(\rho_j-\rho)fg\,dx\right| .
\end{equation*}
$H$ 的单位球在 $C[0,1]$ 中一致有界且一致 $1/2$-H\"older, 故由 Arzel\`a--Ascoli 定理相对紧.
于是乘积族
\begin{equation*}
	\mathcal P=\{fg:\|f\|_H\le1,\|g\|_H\le1\}
\end{equation*}
在 $L^1(0,1)$ 中相对紧. 弱星收敛意味着线性泛函 $\phi\mapsto\int(\rho_j-\rho)\phi$ 在每个固定
$\phi\in L^1$ 上趋于零, 而其范数由 $2R$ 一致控制. 对紧集 $\mathcal P$ 取有限网并使用一致有界性,
收敛便在 $\mathcal P$ 上一致, 得到结论.
\end{proof}

\begin{proposition}[最大值与最小值同时达到]\label{prop:attain}
对每个固定 $n$, 映射 $\rho\mapsto\lambda_n(\rho)$ 在 $A_R$ 的弱星拓扑下连续. 因此 $D_n$ 在
$A_R$ 上达到最大值和最小值.
\end{proposition}

\begin{proof}
由引理 \ref{lem:opnorm} 与紧自伴算子的极小极大原理, 固定编号的特征值 $\mu_k(T_{\rho_j})\to
\mu_k(T_\rho)$. 这些数严格为正, 故取倒数并使用 \eqref{eq:muk} 得 $\lambda_k(\rho_j)\to
\lambda_k(\rho)$, 从而 $D_n$ 弱星连续.

允许集 $A_R$ 在 $L^\infty=(L^1)^*$ 中有界且弱星闭. Banach--Alaoglu 定理给出弱星紧性; 由于
$L^1(0,1)$ 可分, 在该有界集上也可用序列表述. 连续实函数 $D_n$ 因此同时取得最大值与最小值.
\end{proof}

\section{证明义务与解除位置对照表}

\begin{center}
\begin{tabular}{ll}
\toprule
证明义务 & 本文解除位置\\
\midrule
可测正权下的 $C^1$ 迹与简单零点 &
式 \eqref{eq:regularity} 及其后的初值唯一性论证.\\
第 $k$ 模态恰有 $k-1$ 个内部零点 &
Pr\"ufer 角公式 \eqref{eq:theta_x} 与第 2.2 节.\\
相邻模态严格交错且无共同内部零点 &
引理 \ref{lem:interlace}.\\
Wronskian 的符号约定与全区间严格性 &
式 \eqref{eq:wdef}--\eqref{eq:wneg}.\\
首、中、尾结点区间的商极限 &
第 4.1 节的值域表.\\
精确零点数、端点等号、零点简单性 &
命题 \ref{prop:count} 与说明 \ref{rem:n2}.\\
任意 $L^\infty$ 方向的谱隙导数 &
引理 \ref{lem:fh}.\\
最大化和最小化的一侧变分符号 &
命题 \ref{prop:saturation}.\\
零点到开关以及开关到零点的双向蕴含 &
命题 \ref{prop:switchzero}.\\
先得到有限分块、再定义块能量 &
命题 \ref{prop:switchzero} 后的说明与第 6 节首段.\\
接口跳量保留 $r_+-r_-$ 因子 &
式 \eqref{eq:jump}.\\
全局常数中因子 $2$ 的来源 &
式 \eqref{eq:kinetic}--\eqref{eq:kintegral}.\\
端点绝对值不等式升级为有向比值 &
式 \eqref{eq:parity} 与 \eqref{eq:qstrict}.\\
每个极值子而非选定分支 &
定理 \ref{thm:main} 前的量词约定及第 8.2 节.\\
最大值与最小值的达到 &
附录 A 的命题 \ref{prop:attain}.\\
\bottomrule
\end{tabular}
\end{center}

\section{符号速查}

\begin{center}
\begin{tabular}{ll}
\toprule
符号 & 含义\\
\midrule
$A_R$ & 可测权重盒 $1\le\rho\le R$ a.e.\\
$a,b,D$ & $a=\lambda_n$、$b=\lambda_{n+1}$、$D=b-a>0$\\
$u_n,u_{n+1}$ & 按 $\int\rho u_k^2=1$、$u'_k(0)>0$ 归一化的相邻特征函数\\
$c$ & $c=\sqrt{a/b}\in(0,1)$\\
$F$ & 切换函数 $au_n^2-bu_{n+1}^2$\\
$W$ & Wronskian $u'_{n+1}u_n-u_{n+1}u'_n$, 本文约定下严格为负\\
$Q$ & 在每个 $u_n$ 结点区间定义的商 $u_{n+1}/u_n$\\
$q_0,q_1$ & 左、右端斜率比, 见 \eqref{eq:q}\\
$K$ & 两相邻模态的块能量差, 极值子上全局恒等于 $-2D$\\
$Z(F;(0,1))$ & 切换函数的内部零点集合, 恰等于有效开关集合\\
\bottomrule
\end{tabular}
\end{center}

\begin{thebibliography}{4}
\bibitem{teschl} G. Teschl, \emph{Ordinary Differential Equations and Dynamical Systems},
	Graduate Studies in Mathematics 140, American Mathematical Society, 2012.
	\url{https://www.mat.univie.ac.at/~gerald/ftp/book-ode/ode.pdf}
\bibitem{wm79} B. E. Willner and T. J. Mahar, \emph{Extrema of functions of eigenvalues},
	Journal of Mathematical Analysis and Applications 72(2), 730--739, 1979.
	\url{https://doi.org/10.1016/0022-247X(79)90260-9}
\bibitem{kz96} Q. Kong and A. Zettl, \emph{Eigenvalues of Regular Sturm--Liouville Problems},
	Journal of Differential Equations 131(1), 1--19, 1996.
	\url{https://doi.org/10.1006/jdeq.1996.0154}
\bibitem{sun22} H. Sun, \emph{On the minimum eigenvalue gap for vibrating string},
	Journal of Mathematical Analysis and Applications 516(1), 126513, 2022.
	\url{https://doi.org/10.1016/j.jmaa.2022.126513}
\end{thebibliography}

\end{document}
